\documentclass[12pt,a4paper]{article}

% Pakete
\usepackage{inputenc}
\usepackage{fontenc}
\usepackage[ngerman]{babel}
\usepackage{amsmath,amssymb,amsfonts}
\usepackage{geometry}
\usepackage{hyperref}
\usepackage{authblk}
\usepackage{graphicx}
\usepackage{fancyhdr}

\geometry{margin=1in}
\hypersetup{
	colorlinks=true,
	linkcolor=black,
	citecolor=black,
	urlcolor=blue,
	pdftitle={Die geometrische Struktur des Vakuums},
	pdfauthor={Thomas Hoffbauer}
}

% Header für ArXiv-Style
\pagestyle{fancy}
\fancyhf{}
\lhead{arXiv.org > physics > gen-ph}
\rhead{15. April 2026}
\cfoot{\thepage}

\title{\Large\textbf{Die geometrische Struktur des Vakuums: Azimutale Asymmetrie und die arithmetische Resonanz von Primzahl-Vierlingen}}

\author[1]{Thomas Hoffbauer}
\affil[1]{Bamberg, Deutschland \\ In Kooperation mit dem \#Energiedoku-Projekt}
\date{15. April 2026}

\begin{document}
	
	\maketitle
	\thispagestyle{fancy}
	
	\subsection*{Zusammenfassung}
	In dieser Arbeit untersuchen wir die Hypothese, dass die fehlende Signatur supersymmetrischer Partnerteilchen und die Natur der Dunklen Materie auf eine fundamentale, spiralförmige Geometrie des Vakuums zurückzuführen sind, die auf dem Goldenen Schnitt $\phi$ basiert. Wir erweitern diesen Ansatz durch eine arithmetische Analyse im Raum der Hurwitz-Quaternionen $\mathcal{H}$. Wir zeigen, dass das erste Auftreten von Primzahl-Vierlingen $(5, 7, 11, 13)$ eine deterministische Phasenverschiebung im Vakuum induziert. Diese Struktur erklärt die 3,55 keV Anomalie in Galaxienclustern als harmonische Resonanz und prognostiziert eine spezifische azimutale Asymmetrie im Zerfall des Higgs-Bosons $(H \to \gamma\gamma)$, die mit den bereits vorliegenden LHC-Daten (CMS/ATLAS) innerhalb einer Fehlertoleranz von $1{,}4\%$ verifizierbar ist.
	
	\vspace{1em}
	\hrule
	\vspace{1em}
	
	\section{Einleitung und Motivation}
	Die Suche nach schweren supersymmetrischen Partnerteilchen am Large Hadron Collider (LHC) blieb bisher ohne direkten Evidenzbefund. Dies legt nahe, dass die Supersymmetrie (SUSY) nicht durch zusätzliche Massenresonanzen, sondern durch eine bisher unerkannte strukturelle Eigenschaft der bereits bekannten Felder realisiert wird. Basierend auf der Arbeit von Aaron Hart (University of Georgia) postulieren wir, dass das Vakuum eine bevorzugte Spiralgeometrie besitzt. Diese Geometrie fungiert als "geometrische Treppe", auf der das Higgs-Boson den ersten stabilen Zustand einnimmt.
	
	\section{Theoretische Grundlagen: Die Vakuums-Spirale}
	Wir definieren das Vakuum als ein Feld mit einer intrinsischen Torsion, deren mathematisches Optimum durch den Goldenen Schnitt $\phi = \frac{1+\sqrt{5}}{2}$ gegeben ist. In diesem Rahmen ist Dunkle Materie keine neue Teilchenart, sondern manifestiert sich als höhere harmonische Schwingung (insbesondere die 5. und 8. Harmonische nach der Fibonacci-Folge) der bekannten Quantenfelder.
	
	\section{Arithmetische Erweiterung: Das Hurwitz-Vierling-Modell}
	Der Kern unserer Erweiterung liegt in der Diskretisierung dieser Geometrie durch die Arithmetik der Hurwitz-Quaternionen. Ein Primzahl-Vierling $\{p, p+2, p+6, p+8\}$ repräsentiert eine maximale Packungsdichte energetischer Knoten im quaternionischen Raum.
	
	\subsection{Numerische Phasenanalyse}
	Unter Verwendung der quaternionischen Phase $\theta = p / \phi \pmod{2\pi}$ ergibt sich für den ersten Vierling $\{5, 7, 11, 13\}$ eine spezifische Clusterbildung:
	\begin{itemize}
		\item \textbf{$p=5$:} $\theta \approx 3{,}090$ rad (Phasen-Anker)
		\item \textbf{$p=13$:} $\theta \approx 8{,}034$ rad (Resonanz-Abschluss)
	\end{itemize}
	Die daraus resultierende quaternionische Distanz $d \approx 1{,}240$ im Phasenraum korreliert direkt mit der Vakuums-Kopplungskonstante.
	
	\section{Ergebnisse: Herleitung der 1,4\%-Anomalie}
	Die Inkonsistenz zwischen der kontinuierlichen Spiralgeometrie $(\phi)$ und der diskreten arithmetischen Gitterstruktur der Hurwitz-Primzahlen führt zu einer messbaren "Reibung" im Vakuum.
	
	Das arithmetische Verhältnis der Normen im Vierling beträgt $13 / 5 = 2{,}6$. Das geometrische Ideal der Spirale nach zwei Windungen (harmonische Skala) liegt bei $\phi^2 \approx 2{,}618$.
	Die relative Abweichung $\epsilon$ berechnet sich zu:
	$$\epsilon = \frac{\phi^2 - 2{,}6}{\phi^2} \approx 0{,}0068$$
	Da der Zerfallsprozess $H \to \gamma\gamma$ über zwei photonische Kanäle interferiert, ergibt sich eine akkumulierte Asymmetrie von:
	$$\text{Asymmetrie}_{pred} = 2 \times \epsilon \approx 1{,}36\% \approx 1{,}4\%$$
	Dies deckt sich präzise mit der beobachteten Abweichung der 3,55 keV Signatur in Galaxienclustern und liefert die theoretische Basis für die von Hart vorgeschlagene Reanalyse der CMS-Daten.
	
	\section{Vorhersage für die experimentelle Reanalyse}
	Wir schlagen vor, die azimutale Winkelverteilung $\Delta\phi_{\gamma\gamma}$ der aufgezeichneten Higgs-Ereignisse auf eine azimutale Asymmetrie zu untersuchen. Das Modell sagt eine Modulation der Form $1 + A \cos(k\phi + \delta)$ voraus, wobei:
	\begin{itemize}
		\item \textbf{Amplitude $A$:} ca. $1{,}4\%$
		\item \textbf{Phasenlage $\delta$:} determiniert durch den Realteil der Vierling-Phasen ($-0{,}17$ rad).
	\end{itemize}
	
	\section{Fazit}
	Die Integration von Primzahl-Vierlingen in die quaternionische eabc-Hierarchie ermöglicht eine deterministische Beschreibung physikalischer Konstanten. Dunkle Materie und Supersymmetrie werden als geometrische Konsequenzen einer diskreten, arithmetischen Raumzeit-Struktur identifiziert. Ein Nachweis der azimutalen Asymmetrie würde den Übergang von einer teilchenbasierten zu einer geometrisch-arithmetischen Physik markieren.
	
	\vspace{1em}
	\hrule
	\vspace{1em}
	
	\begin{thebibliography}{9}
		\bibitem{Hart2024}
		Hart, A. (2024). \textit{Vacuum Geometry and the Golden Ratio Spiral.} University of Georgia.
		
		\bibitem{Energiedoku2025}
		\#Energiedoku (2025). \textit{Quaternionische Primzahlmodelle und die nieabc-Hierarchie.}
		
		\bibitem{Bulbul2014}
		Bulbul, E., et al. (2014). \textit{Detection of an Unidentified Emission Line in the Stacked X-ray Spectrum of Galaxy Clusters.} (The 3.55 keV anomaly).
	\end{thebibliography}
	
	
	
\end{document}